\documentclass{article}
\usepackage{graphicx} % Required for inserting images
\usepackage{bm}
\usepackage{amsmath}
\usepackage{csquotes}
\usepackage{xcolor}
\usepackage{parskip}


\usepackage{listings}
\lstdefinestyle{myRstyle}{
  language=R,
  basicstyle=\ttfamily\small,
  commentstyle=\color{gray}\itshape,
  keywordstyle=\color{blue}\bfseries,
  stringstyle=\color{orange},
  numberstyle=\tiny\color{gray},
  numbers=left,
  stepnumber=1,
  frame=single,
  breaklines=true
}

%\usepackage{authblk}
\newcommand{\I}{\mathcal{I}}
\newcommand{\X}{\mathbf{X}}
\newcommand{\tp}{^{\mathrm T}}
\DeclareMathOperator{\tr}{tr}
\author{Da Hoang \\ \texttt{huuda@ucsb.edu}\and Patrick McDonough\\\texttt{pjmcdonough@ucsb.edu}\and Anthony Meena\\\texttt{anthonymeena@ucsb.edu}\and Jordan Wingenroth\\\texttt{jmwingenroth@ucsb.edu}}
\title{Linear Algebra Problem Set \#1}
\date{\today}

\begin{document}

\maketitle

%\section{Linear Algebra 1}
\subsection*{Q1}

\subsubsection*{1(a)}
\[
(AB)\tp = \left( \begin{bmatrix} 2 & 0 \\ 3 & 8 \end{bmatrix} \begin{bmatrix} 7 & 2 \\ 6 & 3 \end{bmatrix} \right)\tp
= \left( \begin{bmatrix} 14+0 & 4+0 \\ 21+48 & 6+24 \end{bmatrix} \right)\tp
\]

\[
= \left( \begin{bmatrix} 14 & 4 \\ 69 & 30 \end{bmatrix} \right)\tp
= \begin{bmatrix} 14 & 69 \\ 4 & 30 \end{bmatrix}
\]

\[
B\tp A\tp = \begin{bmatrix} 7 & 2 \\ 6 & 3 \end{bmatrix}\tp \begin{bmatrix} 2 & 0 \\ 3 & 8 \end{bmatrix}\tp
= \begin{bmatrix} 7 & 6 \\ 2 & 3 \end{bmatrix} \begin{bmatrix} 2 & 3 \\ 0 & 8 \end{bmatrix}
= \begin{bmatrix} 14+0 & 21+48 \\ 4+0 & 6+24 \end{bmatrix}
\]

\[
= \begin{bmatrix} 14 & 69 \\ 4 & 30 \end{bmatrix}
\]

\textbf{1(b)}
\[
(AB)^{-1} = \begin{bmatrix} 14 & 4 \\ 69 & 30 \end{bmatrix}^{-1}
= \frac{1}{144} \begin{bmatrix} 30 & -4 \\ -69 & 14 \end{bmatrix}
\]

\[
B^{-1}A^{-1} = \frac{1}{9} \begin{bmatrix} 3 & -2 \\ -6 & 7 \end{bmatrix} \cdot \frac{1}{16} \begin{bmatrix} 8 & 0 \\ -3 & 2 \end{bmatrix}
= \frac{1}{144} \begin{bmatrix} 30 & -4 \\ -69 & 14 \end{bmatrix}
\]

\textbf{1(c)}
\[
(A\tp)^{-1} = \left( \begin{bmatrix} 2 & 0 \\ 3 & 8 \end{bmatrix}\tp \right)^{-1}
= \left( \begin{bmatrix} 2 & 3 \\ 0 & 8 \end{bmatrix} \right)^{-1}
= \frac{1}{16} \begin{bmatrix} 8 & -3 \\ 0 & 2 \end{bmatrix}
\]

\[
(A^{-1})\tp = \left( \begin{bmatrix} 2 & 0 \\ 3 & 8 \end{bmatrix}^{-1} \right)\tp
= \left( \frac{1}{16} \begin{bmatrix} 8 & 0 \\ -3 & 2 \end{bmatrix} \right)\tp
= \frac{1}{16} \begin{bmatrix} 8 & -3 \\ 0 & 2 \end{bmatrix}
\]



\vfill

\subsection*{Q2}
This is not the case, in general.
Consider the counterexample \[\det(\I_{n})+\det(\I_{n})=2\ne 2^n=\det(\I_{n}+\I_{n})\]
\subsection*{Q3}
\subsubsection*{3(a)} The Hessian is 
\[
 \begin{bmatrix} -2 & 1 \\ 1 & -2 \end{bmatrix}
\] which has principal minors $-2$ and $3$. Therefore it's negative definite and therefore it's strictly concave.


\subsubsection*{3(b)} The Hessian is 
\[
 \begin{bmatrix} 6 & -2 & 0 \\  -2 & 4 & 2 \\ 0 & 2 & 2  \end{bmatrix}
\]
It's easy enough to see that its first two principal minors are $6$ and $20$. The determinant is $16$.
Therefore, the matrix is positive definite and the function is strictly convex.

\subsection*{Q4}

\subsubsection*{Symmetry}

To show that $M_{\X}$ is symmetric,
\begin{align*}
    \mathcal M_{\X}\tp&=(\I_n-\X(\X\tp\X)^{-1}\X\tp)\tp\\
    &=\I_n\tp-(\X(\X\tp\X)^{-1}\X\tp)\tp&\text{by the distributive property}\\
    &=\I_n-(\X\tp)\tp(\X(\X\tp\X)^{-1})\tp&\text{by the reverse-order law and }\I_n\tp=\I_n\\
    &=\I_n-\X((\X\tp\X)^{-1})\tp\X\tp&\text{by the reverse-order law and }(A\tp)\tp=A\\
    &=\I_n-\X((\X\tp\X)\tp)^{-1}\X\tp&\text{since the inverse and transpose commute}\\
    &=\I_n-\X(\X\tp\X)^{-1}\X\tp&\text{by }(\X\tp\X)\tp=\X\tp\X\\
    &=\mathcal M_{\X}
\end{align*}
Because $\mathcal M_{\X}\tp=\mathcal M_{\X}$, the matrix $\mathcal M_{\X}$ is symmetric.

\subsubsection*{Idempotency}

To show that $\mathcal M_{\X}$ is idempotent, write it as a difference of two matrices,
\begin{equation*}
    \mathcal M_{\X}=\mathbf A-\mathbf B,\qquad\text{where }\mathbf A=\I_n\text{ and }\mathbf B=\X(\X\tp\X)^{-1}\X\tp.
\end{equation*}
Each of the four products appearing in $(\mathbf A-\mathbf B)^2$ simplifies:
\begin{align*}
    \mathbf A^2&=\I_n\I_n=\I_n=\mathbf A\\
    \mathbf A\mathbf B&=\I_n\X(\X\tp\X)^{-1}\X\tp=\X(\X\tp\X)^{-1}\X\tp=\mathbf B\\
    \mathbf B\mathbf A&=\X(\X\tp\X)^{-1}\X\tp\I_n=\X(\X\tp\X)^{-1}\X\tp=\mathbf B\\
    \mathbf B^2&=\X(\X\tp\X)^{-1}\underbrace{\X\tp\X(\X\tp\X)^{-1}}_{=\,\I}\X\tp=\X(\X\tp\X)^{-1}\X\tp=\mathbf B
\end{align*}
Expanding the square and substituting these four results gives
\begin{align*}
    \mathcal M_{\X}^2&=(\mathbf A-\mathbf B)^2\\
    &=\mathbf A^2-\mathbf A\mathbf B-\mathbf B\mathbf A+\mathbf B^2\\
    &=\mathbf A-\mathbf B-\mathbf B+\mathbf B\\
    &=\mathbf A-\mathbf B\\
    &=\mathcal M_{\X}
\end{align*}
so $\mathcal M_{\X}$ is idempotent.

\subsection*{Q5}
The characteristic polynomial is 
\begin{align*}
    (.8-\lambda)(.95-\lambda)-(.05)(.2)&=.76-1.75\lambda+\lambda^2-.01\\
    &=\lambda^2-1.75\lambda+.75\\
    &=\frac14(4\lambda^2-7\lambda+3)\\
    &=\frac14(\lambda-1)(4\lambda-3)
\end{align*}
Therefore, the eigenvalues are $\lambda_1=1$ and $\lambda_2=\frac{3}{4}$

For $\lambda_1=1$:
\begin{align*}
 (C - \lambda_1 I) u & = 0 \\
\begin{bmatrix} 0.8-1 & 0.05 \\ 0.2 & 0.95-1 \end{bmatrix} \begin{bmatrix} u_1 \\ u_2 \end{bmatrix} &= 0 \\
-0.2u_1 + 0.05u_2 &= 0 \\
u_2 &= 4u_1
\end{align*}
Normalizing the vector, we get the eigenvector for $\lambda_1=1$ is: $u = \begin{bmatrix} \dfrac{1}{\sqrt{17}} & \dfrac{4}{\sqrt{17}} \end{bmatrix}$

For $\lambda_2=\frac{3}{4}=0.75$:
\begin{align*}
 (C - \lambda_2 I) u & = 0 \\
\begin{bmatrix} 0.8-0.75 & 0.05 \\ 0.2 & 0.95-0.75 \end{bmatrix}  \begin{bmatrix} u_1 \\ u_2 \end{bmatrix} &= 0 \\
0.05u_1 + 0.05u_2 &= 0 \\
u_2 &= -u_1
\end{align*}
Normalizing the vector, we get the eigenvector for $\lambda_2=3/4$ is: $u = \begin{bmatrix} \dfrac{1}{\sqrt{2}} & \dfrac{-1}{\sqrt{2}} \end{bmatrix}$

\subsection*{Q6}
The sum given in (a) is $\mathbf u'(\frac12\Sigma^{-1})\mathbf u$ and the one given in (b) is $\mathbf u'\Lambda\mathbf u$.

\subsection*{Q7}
The identities given in (a) follow directly from the definition of matrix multiplication.

The identity in (b) results from applying the identities in (a) to the identity $(\X\tp\X)^{-1}X\tp \mathbf y=(\frac1n\X\tp\X)^{-1}\frac1n\X\tp \mathbf y$ which follows from the commutativity and associativity of scalar-matrix multiplication and the identity $(\frac1n\X)^{-1}=n\X^{-1}$ which can be obtained by applying the commutativity and associativity of scalar-matrix multiplication to show that $(\frac1n\X)^{-1}n\X^{-1}=\I$. 

(c) We have:
\begin{align*}
    \hat{\beta} & = (X\tp X)^{-1} X\tp Y \\
    \hat{\beta} & = (X\tp X)^{-1} X\tp (X \beta + u) \\
    \hat{\beta} & = (X\tp X)^{-1} X\tp X \beta + (X\tp X)^{-1} X\tp u \\
    \hat{\beta} & = I \beta + (X\tp X)^{-1} X\tp u \\
    \hat{\beta} - \beta & = (X\tp X)^{-1} X\tp u \\
    \sqrt{n} (\hat{\beta} - \beta) & = \sqrt{n} (\frac{1}{n} X\tp X)^{-1} (\frac{1}{n} X\tp u) \text{ - this is because } (\frac{1}{n})^{-1} = n
\end{align*}

We already had that $(X\tp X)^{-1} = (\sum_{i=1}^{n} x_i x_i\tp)^{-1}$, and we also have $X\tp u = \sum_{i=1}^{n} x_i u_i$. Therefore, the formula is proved.

(d) We have that $\hat{u} = \mathcal M_{\X} \ y = \mathcal M_{\X} \ (X\beta + u) = \mathcal M_{\X} X \beta + \mathcal M_{\X} u $.

First, we want to prove that $\mathcal M_{\X} X \beta = 0$.
\begin{align*}
    \mathcal M_{\X} X \beta & = [I - X (X\tp X)^{-1} X\tp] X \beta \\
    \mathcal M_{\X} X \beta & = X \beta - X (X\tp X)^{-1} X\tp X \beta \\
    \mathcal M_{\X} X \beta & = X \beta - X I \beta \\
    \mathcal M_{\X} X \beta & = 0
\end{align*}

Therefore, $\hat{u} = \mathcal M_{\X} X \beta + \mathcal M_{\X} u = \mathcal M_{\X} u $.

Hence, $\sum (\hat{u_i})^2 = \hat{u}\tp \hat{u} = (\mathcal M_{\X} u)\tp (\mathcal M_{\X} u) = u\tp (\mathcal M_{\X})\tp \mathcal M_{\X} u = u\tp \mathcal M_{\X} u $.

The last step is from the fact that $\mathcal M_{\X}$ is symmetric, therefore, ($\mathcal M_{\X})\tp = \mathcal M_{\X}$. $\mathcal M_{\X}$ is also idempotent, so $\mathcal M_{\X} \mathcal M_{\X} = \mathcal M_{\X}$.

\pagebreak
\subsection*{Q8}

Two functions that satisfy the specifications are below:



\begin{lstlisting}[style=myRstyle]
### This function solve a linear system of equations: Ax = b
# A: nxn square matrix; b: nx1 vector
# If A is singular, print out "It does not have a unique solution"
# If A is not singular, output the solution 

MC_solve = function (A, b) {
  
  ### Calculate det_A which is the determinant of A
  #Note: round to 10 decimal place, because even if A is singular, 
  #det(A) may be a very small number not equal to 0 (because of R's algorithm)
  det_A = round(det(A), 10)
  
  ### Check if det_A is equal to 0, if not, solve the equation
  if (det_A == 0){
    print("It does not have a unique solution")
  } else {
    solution = solve(A, b)
    solution
  }
}

### For problem (a)
A1 = matrix(c(1, -2, 1, 0, 2, -8, -4, 5, 9), nrow = 3, byrow = TRUE)
b1 = c(0, 8, -9)

# Call the function on A1, b1
MC_solve(A1, b1)
### Result: (29, 16, 3)

### For problem (b)
A2 = matrix(c(1, 4, 2, 2, 5, 1, 3, 6, 0), nrow = 3, byrow = TRUE)
b2 = c(1, -2, 3)

# Call the function on A2, b2
MC_solve(A2, b2)
### Result: print "It does not have a unique solution"
\end{lstlisting}

\end{document}
